% Part: computability
% Chapter: computability-theory
% Section: coding-computations

\documentclass[../../../include/open-logic-section]{subfiles}

\begin{document}

\olfileid[pt]{cmp}{thy}{cod}
\olsection{Codificação de computações}

Em todo modelo de computação, é possível fazer o seguinte:
\begin{enumerate}
\item Descrever as \emph{definições} de funções computáveis de maneira
  sistemática. Por exemplo, pode-se pensar que especificações de
  máquinas de Turing, definições recursivas ou programas em uma
  linguagem de programação fornecem essas definições.
\item Descrever o registro completo da computação de uma função dada
  por alguma definição para uma determinada entrada. Por exemplo, uma
  computação de máquina de Turing pode ser descrita pela sequência de
  configurações (estado da máquina, conteúdo da fita) correspondente
  a cada passo da computação.
\item Testar se um suposto registro de uma computação é de fato o
  registro de como uma função computável com uma dada definição seria
  calculada para uma determinada entrada (na qual a função está
  definida, isto é, a computação para).
\item Extrair de tal descrição do registro completo de uma computação
  o valor da função para uma determinada entrada. Por exemplo, o
  conteúdo da fita no último passo de uma computação de máquina de
  Turing que para é o valor.
\end{enumerate}

Por meio de codificação, é possível atribuir a cada descrição de uma
função computável um \emph{índice} numérico de tal maneira que as
instruções possam ser recuperadas computavelmente a partir do índice.
De modo semelhante, o registro completo de uma computação também pode
ser codificado por um único número. A relação aritmética resultante
``$s$~codifica o registro da computação da função de índice~$e$ para
a entrada~$x$'' e a função ``saída da sequência de computação de
código~$s$'' são, então, computáveis; na verdade, são recursivas
primitivas.

Esse fato fundamental é muito poderoso e nos permite demonstrar
diversos resultados notáveis e importantes sobre a computabilidade,
independentemente do modelo de computação escolhido.

\end{document}

